% !TEX root = ../metatime_monograph.tex
\chapter{Information-Spinor Crosswalk and Dependency Direction}
\label{app:information_spinor_crosswalk}

% CITATION_CONTEXT_V10_9_BEGIN
\paragraph{Established context.}
Binary information, projective quantum-state geometry, and geometric holonomy
follow the standard Shannon, Simon, Berry, and differential-geometric context
\cite{shannon1948,Simon1983,berry1984,BottTu1982}.  The identity-axis,
cycle-count, sector pairing, and CIEL/PhaseNav dependency assignments below are
specific to the TIR/Metatime programme.
% CITATION_CONTEXT_V10_9_END

\claimstatus{Mixed: exact mathematical identities plus explicitly marked Metatime/TIR model assignments.}

\section{Why this crosswalk exists}
The dedicated critical-axis programme under
\path{TIR/zeta_information_axis/}
formalizes the repeated appearance of the balanced coordinate $1/2$ in binary information, complement symmetry, Bloch geometry, Berry holonomy, destructive interference, centred reciprocal geometry, and the anti-linear zeta involution.  The parent Metatime monograph uses some of the same structures in a broader phenomenological programme.  This appendix records the direction of dependence so that a result from one layer is not silently promoted in the other.

The governing rule is:
\[
\boxed{
\text{exact mathematical bridge}
\;\neq\;
\text{physical identification}
\;\neq\;
\text{empirical validation}.
}
\]

\section{The half axis as an identity coordinate}
The binary complement
\[
C(\sigma)=1-\sigma
\]
has the unique fixed point
\[
\sigma=\frac12.
\]
Using the centered coordinate
\[
x=\sigma-\frac12,
\]
complement becomes the orientation reversal $x\mapsto-x$.  Hence negative, zero, and positive centered values can be read geometrically as the two orientations and the fixed axis of the same one-dimensional chart.  The same centering idea is used by the critical-axis submonograph for $u=s-1/2$.

In the qubit representation,
\[
n_z=2\sigma-1,
\]
so $\sigma=1/2$ is the Bloch equator.  Phase evolution can remain nontrivial while the identity coordinate remains balanced.  This is the sense in which the half value is an \emph{axis} rather than a frozen dynamical state.

\section{Normalized cycles before radians}
The critical-axis solver introduces a turn coordinate
\[
q\in\mathbb R/\mathbb Z
\]
with intrinsic winding frequency
\[
f_q=\frac{\Delta N}{\Delta t}.
\]
An angular representation is introduced only afterwards by
\[
\phi=Cq.
\]
For radians, $C=2\pi$.  The sequence
\[
\text{recurrence}\to\text{winding}\to\text{angle}
\]
is therefore conceptually prior to the conventional equation $\omega=2\pi f$.

This cross-references the phase and Poincar\'e structures of Chapter~\ref{ch:metatime_framework} and Appendix~\ref{app:poincare}, while retaining the publication claim boundary of Appendix~\ref{app:kappa}.

\section{The $24$ cross-factorization}
The exact integer identity is
\begin{equation}
24=8\cdot3=12\cdot2=6\cdot4.
\label{eq:appP-24}
\end{equation}
The equality is arithmetic.  Interpretations of these factors are model-level:
\begin{itemize}
\item $8\cdot3$: eight declared mixing sectors times three flavours;
\item $12\cdot2$: twelve projective cycles times two half-turn units;
\item $6\cdot4$: six spinor cycles times four half-turn units.
\end{itemize}
The three-flavour layer connects structurally to Chapter~\ref{ch:pmns_mixing}.  The projective/spinor interpretation connects to the Berry and holonomy structures reviewed in Chapter~\ref{ch:metatime_framework}.  None of these labels changes the status of $24$ from a discrete model normalization into an independently established physical constant.

\section{Arithmetic--geometry factorization of $\kappa$}
If the Metatime model assigns one binary information quantum to twelve normalized projective cycles,
\begin{equation}
\frac{\dd I}{\dd q}=\frac{\ln2}{12},
\label{eq:appP-info-turn}
\end{equation}
and if the chosen angular representation has closure period $C$, then
\begin{equation}
\frac{\dd I}{\dd\phi}=\frac{\ln2}{12C}.
\end{equation}
For $C=2\pi$ this reproduces
\begin{equation}
\kappa=\frac{\ln2}{24\pi}.
\label{eq:appP-kappa}
\end{equation}
Equation~\eqref{eq:appP-kappa} is exact conditional arithmetic once the cycle assignment and angular closure have been declared.  The independent physical derivation of the cycle assignment remains a model debt, consistently with Appendix~\ref{app:kappa}.

\section{Paired structures and implication directions}
\begin{longtable}{p{0.26\textwidth}p{0.29\textwidth}p{0.33\textwidth}}
\toprule
Source structure & Paired consequence & Status / direction \\
\midrule
$\sigma=1/2$ & $C(\sigma)=\sigma$ & Exact both-way characterization in the binary interval.\\
$\sigma=1/2$ & $H_{\rm bin}=\ln2$ & Exact; entropy maximum uniquely selects the half point.\\
$\sigma=1/2$ + qubit representation & Bloch equator $n_z=0$ & Exact inside the stated two-level representation.\\
Bloch equator + one azimuthal loop & Berry holonomy $-1$ & Standard geometric-phase result in the chosen state family.\\
Symmetric complementary readout + half-turn phase & exact destructive cancellation & Exact; the converse also selects $\sigma=1/2$.\\
$u=s-1/2$ + $J(s)=1-\bar s$ & fixed axis $\Re s=1/2$ & Exact/classical.\\
Centered zeta chart + $u\mapsto1/u$ & critical-axis set invariance & Exact; does not imply zero-set invariance.\\
Projective cycle + spin $1/2$ & double-cover period & Standard representation theory.\\
$\ln2/12$ per projective turn + $C=2\pi$ & $\kappa=\ln2/(24\pi)$ & Conditional arithmetic; cycle assignment is model-level.\\
TIR sector-holonomy operator & mass-spectrum candidate & Retrospective model layer; see Appendix~\ref{app:sector_holonomy_release}.\\
Critical-axis cancellation representation of every zeta zero & RH & Open bridge; not derived by Metatime holonomy alone.\\
\bottomrule
\end{longtable}

\section{Holonomy policy}
A relation edge is not automatically a holonomy edge.  Holonomy is meaningful only when a connection and a closed or composable path are defined.  This distinction matters for the parent programme because ``paired'' structures may be algebraic, geometric, semantic, or empirical without carrying a phase.

The CIEL/PhaseNav runtime therefore stores a normalized holonomy coordinate only on eligible path edges.  Generated PhaseNav geometry remains a candidate until independent TIR provenance supplies the semantic relation.  Conversely, a TIR relation without an executable geometric path does not acquire a synthetic phase merely to complete a diagram.

\section{Runtime dependency direction}
The companion CIEL/PhaseNav package is organized so that higher-level operators depend on primitive geometry rather than the reverse:
\[
\boxed{
\begin{array}{c}
\text{identity axis / centered coordinate}\\
\text{normalized cycle }\mathbb R/\mathbb Z\\
\text{phase crystal }(\mathbb R/\mathbb Z)^{36}\\
\text{typed relation edges}
\end{array}
\longrightarrow
\text{tangent/log-map vectorization}
\longrightarrow
\text{semantic and holonomic solvers}.
}
\]
On the phase crystal, the shortest oriented tangent displacement naturally lies in $[-1/2,1/2)$ per normalized coordinate.  Its sign is therefore generated by orientation; an identical coordinate produces zero.  This is the concrete computational meaning of the statement that signed vector components emerge from the geometry rather than being attached afterwards as semantic labels.

\section{Claim firewall}
The dedicated solver uses four statuses: \texttt{EXACT}, \texttt{STANDARD}, \texttt{MODEL}, and \texttt{OPEN}.  Exact closure permits only the first two classes.  Model closure may additionally use model assignments.  Open rules are never automatically promoted.

In particular, the solver can reproduce the mutually consistent half-axis routes and the conditional normalization \eqref{eq:appP-kappa}, while still reporting the canonical zeta zero-state representation as a missing premise.  This preserves the scientific-status policy of this publication candidate: technical consistency and structural unification do not by themselves constitute physical or number-theoretic closure.
